(1) 在行星静止参考系，探测器以相对速度 \(u'=u-V<0\) 迎向行星；理想一维弹性折返令 \(u'_{\rm out}=-(u-V)=V-u\)。因此回到实验室系
\[
u_{\rm out}=V+(V-u)=2V-u.
\]
在 V>u 条件下，探测器的末速度大于 V，得到向前的助推。

(2) 动能差为
\[
\Delta K=\frac m2[(2V-u)^2-u^2]=2mV(V-u)>0.
\]
若选择相反的几何关系 V<u，上式改为负数，说明在一维模型中并非任何行星掠过都必然使探测器增能。

(3) 当 M 有限，令末速度为 \(u_f,V_f\)。动量守恒 \(mu+MV=mu_f+MV_f\)，能量守恒 \(mu^2+MV^2=mu_f^2+MV_f^2\)，排除无散射平凡支后求得
\[
u_f=\frac{m-M}{m+M}u+\frac{2M}{m+M}V,\qquad
V_f=\frac{2m}{m+M}u+\frac{M-m}{m+M}V.
\]
将两式代回总动量和动能可独立验证守恒。

(4) 直接相减得 \(u_f-V_f=-(u-V)\)，即弹性碰撞的相对速度反号。当 M≫m，\(u_f\to2V-u\)，而行星的速度变化 \(V_f-V=2m(u-V)/(m+M)<0\)。行星的巨大总动能因很小的速度变化而减少，供给探测器增益；若 V=0，则探测器仅反转方向而不增加实验室动能。
